\section{Preliminaries}

In this section we recall the material we will use in the sequel.
In the f\/irst subsection we recall the coalgebra-Galois extensions, as well as the algebra-Galois coextensions.
In the second subsection we discuss the Pontryagin like duality between discrete and linearly compact vector spaces.
Upon recalling the Takeuchi--Galois correspondence between left comodule subalgebras and right module quotient
coalgebras of a~Hopf algebra, we show Pontryagin self-duality of this correspondence.
In the third subsection we recall the relative cyclic homology of algebra extensions.
Finally, the Hopf-cyclic homology, with coef\/f\/icients, of coalgebras is recalled in the fourth subsection.

Throughout the paper, all algebras, coalgebras, and Hopf algebras are over a~f\/ield~$k$, and similarly all unadorned
tensor product symbols $\otimes$ are also over~$k$.
Suppressing the summation, we will write $\Delta(c)=c\ps{1}\otimes c\ps{2}$ for a~comultiplication,
$\rho(v)=v\ns{0}\otimes v\ns{1}$ for a~right coaction, and $\lambda(v)=v\ns{-1}\otimes v\ns{0}$ for a~left coaction.

\subsection{Coalgebra-Galois extensions and algebra-Galois coextensions}

In this subsection we recall the def\/inition and basic properties of the coalgebra-Galois extensions, and the
algebra-Galois coextensions from~\cite{BrzeHaja99,BrzeHaja09}.
Fix the base f\/ield~$k$.

We begin with the coalgebra-Galois extensions.
Let ${\cal C}$ be a~coalgebra coaugmented by the choice of a~group-like element $e\in {\cal C}$.
We say that an algebra ${\cal A}$ is a~\emph{${\cal C}$-algebra} if there is given a~right entwining $\psi\colon  {\cal
C}\otimes {\cal A}\rightarrow {\cal A}\otimes {\cal C}$.
Then ${\cal A}$ becomes a~right ${\cal C}$-comodule via $\rho\colon {\cal A} \to {\cal A} \otimes {\cal C}$ def\/ined as
$\rho(a) =a\ns{0} \otimes a\ns{1}:=\psi(e\otimes a)$.
Then the subspace ${\cal A}^{\co {\cal C}}$ of \emph{coaction invariants} of ${\cal A}$ def\/ined as
\begin{gather*}
{\cal A}^{\co {\cal C}}:=\{b \in {\cal A} \mid \psi(e\otimes b)=b\otimes e\}
\end{gather*}
is a~subalgebra of ${\cal A}$.
Note that ${\cal A}^{\co {\cal C}}$ can be def\/ined as a~limit of a~diagram in vector spaces
\begin{gather*}
{\cal A}^{\co{\cal C}}={\rm Eq}({\cal A}\rightrightarrows {\cal A}\otimes {\cal C})
\end{gather*}
consisting of a~pair of maps $a\mapsto a\ns{0} \otimes a\ns{1}$ and $a\mapsto a\otimes e$.

An algebra extension ${\cal B}\to{\cal A}$ is called a~\emph{${\cal C}$-extension} if ${\cal A}$ is a~${\cal C}$-algebra
and ${\cal B} ={\cal A}^{\co {\cal C}}$.
Finally, a~${\cal C}$-extension ${\cal B}\to{\cal A}$ is said to be \emph{Galois} if the left ${\cal A}$-linear right
${\cal C}$-colinear map
\begin{gather}
\label{aux-canonical-map}
\operatorname{can}\colon \ {\cal A} \otimes_{\cal B} {\cal A} \longrightarrow {\cal A} \otimes {\cal C},
\qquad
a\otimes_{\cal B} a' \mapsto aa'\ns{0} \otimes a'\ns{1}
\end{gather}
is bijective.
The map~\eqref{aux-canonical-map} is called the \emph{canonical map} of the ${\cal C}$-extension.

A \emph{standard} example of a~${\cal C}$-algebra ${\cal A}$ comes from any coaugmented right ${\cal H}$-module
coalgebra~${\cal C}$ and any right ${\cal H}$-comodule algebra ${\cal A}$ over a~Hopf algebra ${\cal H}$.
If we denote by $a\mapsto a_{[0]}\otimes a_{[1]}$ the right ${\cal H}$-coaction on ${\cal A}$ the corresponding
entwining reads as $\psi(c\otimes a)=a_{[0]}\otimes c\cdot a_{[1]}$.

An interesting, due to the geometric examples it covers, class of standard coalgebra-Galois extensions consists of so
called \emph{quotient coalgebra-Galois extensions} def\/ined as follows.
In this setting, one lets~${\cal H}$ to be a~Hopf algebra, a~coaugmented right ${\cal H}$-module coalgebra ${\cal C}$
a~quotient~${\cal H}/{\cal I}$ of~${\cal H}$ by a~coideal right ideal ${\cal I} \subseteq {\cal H}$, and ${\cal A}$
a~right ${\cal H}$-comodule algebra via $\widetilde{\rho}\colon {\cal A} \longrightarrow {\cal A} \otimes {\cal H}$.
Then the composition
\begin{gather}
\label{rho}
\xymatrix{
\rho\colon \ {\cal A} \ar[r]^-{\widetilde{\rho}}&{\cal A} \otimes {\cal H} \ar[r]^-{{\cal A} \otimes\pi}&{\cal A} \otimes {\cal H}/{\cal I}}
\end{gather}
expresses the standard right ${\cal H}/{\cal I}$-coaction on ${\cal A}$.
Therefore the subalgebra of the coaction invariants can be expressed diagrammatically in vector spaces as the equalizer
\begin{gather}
\label{diageq}
{\cal A}^{\co{\cal H}/{\cal I}}={\rm Eq}({\cal A}\rightrightarrows {\cal A}\otimes{\cal H}/{\cal I}),
\end{gather}
where one arrow is~$\rho$ def\/ined by~\eqref{rho}, and the other is the composition
\begin{gather*}%\label{anoar}
\xymatrix{{\cal A} \ar[r]^-{\cong}&{\cal A} \otimes k\ar[r]^-{{\cal A} \otimes\eta}&{\cal A} \otimes {\cal
H}\ar[r]^-{{\cal A} \otimes\pi} &  {\cal A} \otimes {\cal H}/{\cal I}.}
\end{gather*}
Finally, a~Galois ${\cal H}/{\cal I}$-extension ${\cal A}^{\co {\cal H}/{\cal I}} \to {\cal A}$ is called
a~\emph{quotient coalgebra-Galois extension}.
If~${\cal I}$ is a~Hopf ideal, we call such an extension \emph{quotient Hopf--Galois extension}.

The quantum instanton bundle of~\cite{BoneCiccDabrTarl04,BoneCiccTarl02} is a~quotient coalgebra-Galois extension, and
it uses the full generality of the quotient coalgebra-Galois extensions as ${\cal A} \neq {\cal H}$ and ${\cal I} \neq 0$.

For ${\cal I} =0$, the quotient coalgebra-Galois extensions recover the Hopf--Galois extensions, and in case of ${\cal
A} ={\cal H}$ they are called \emph{homogeneous coalgebra-Galois extensions}.
In particular, viewing ${\cal H}$ as a~right ${\cal H}$-comodule algebra via its comultiplication, one obtains the
homogeneous ${\cal H}/{\cal I}$-Galois extensions
\begin{gather*}
\xymatrix{
\rho\colon \ {\cal H} \ar[r]^{\Delta}& {\cal H} \otimes {\cal H} \ar[r]^{{\cal H} \otimes \pi\ \ \ \ \ \ \ \ \ \ \ \ \ \ } &
{\cal H}\otimes {\cal H}/{\cal I}, \qquad  h \mapsto h\ps{1} \otimes \wbar{h\ps{2}}.
}
%\xymatrix{{\rho\colon \  {\cal H}}   \ar[r]^{\Delta} & {\cal H} \otimes {\cal H}
%\ar[r]^{{\cal H} \otimes\pi}  & {{\cal H} \otimes {\cal H}/{\cal I}},
%\qquad h \mapsto h\ps{1} \otimes \wbar{h\ps{2}}.  }
\end{gather*}
In the context of the faithfully f\/lat homogeneous extensions, the diagrammatical def\/inition of invariants of the
coaction~\eqref{diageq} plays an important role in the context of Pontryagin self-duality of the Takeuchi--Galois
correspondence considered in the next subsection.

Such Galois extensions correspond to the quantum homogeneous spaces considering ${\cal H}={\cal O}(G)$ the algebra of
functions on a~quantum group $G={\rm Spec}({\cal H})$, and ${\cal B}={\cal O}(X)$ the algebra of functions on a~quantum
space $X={\rm Spec}({\cal B})$.
Then the~$G$-action on~$X$ is encoded by the ${\cal H}$-coaction~\cite[Section~1]{DijkKoor94}.
There is a~celebrated example of this construction due to Podle\'s~\cite{Brze96,Podl87} which is a~quantum spherical
f\/ibration $SU_q(2)\rightarrow S^2_{q,\mu, \nu}$, see~\cite{Brze97}, a~quantum deformation of the classical Hopf
f\/ibration of a~3-sphere over a~2-sphere into circles.

We now recall the algebra-Galois coextensions.
Let ${\cal B}$ be an algebra augmented by the choice of a~character $y\colon  {\cal B}\rightarrow k$.
We say that a~coalgebra ${\cal D}$ is a~\emph{${\cal B}$-coalgebra} if there is given a~left entwining $\varphi\colon {\cal
B}\otimes {\cal D}\rightarrow {\cal D}\otimes {\cal B}$.
Then ${\cal D}$ becomes a~left ${\cal B}$-module via $\ell\colon {\cal B} \otimes {\cal D} \to {\cal D}$ def\/ined as
$\ell=({\cal D}\otimes y)\circ \varphi$.
Then the quotient space ${\cal D}/{\cal B}^{+}{\cal D}$ of \emph{action coinvariants} of ${\cal D}$ def\/ined~by
\begin{gather*}
{\cal B}^{+}:=\{b \in {\cal B} \mid y(b)=0\}
\end{gather*}
is a~quotient coalgebra of ${\cal D}$.
Note that ${\cal D}/{\cal B}^{+}{\cal D}$ can be def\/ined as a~colimit of a~diagram in vector spaces
\begin{gather*}
{\cal D}/{\cal B}^{+}{\cal D}={\rm Coeq}({\cal D}\leftleftarrows {\cal B}\otimes {\cal D})
\end{gather*}
consisting of a~pair of maps $\ell\colon b\otimes d\mapsto b\cdot d$ and $b\otimes d\mapsto y(b)d$.

A coalgebra coextension ${\cal D}\to{\cal C}$ is called a~\emph{${\cal B}$-coextension} if ${\cal D}$ is a~${\cal
B}$-coalgebra and ${\cal C} ={\cal D}/{\cal B}^{+}{\cal D}$.
Finally, a~${\cal B}$-coextension ${\cal D}\to{\cal C}$ is said to be \emph{Galois} if the left ${\cal B}$-module, right
${\cal D}$-comodule map
\begin{gather}
\label{aux-cocan}
\operatorname{cocan}\colon \ {\cal B}\otimes {\cal D}\longrightarrow {\cal D}\Box_{\cal C}{\cal D},
\qquad
b\otimes d\mapsto b\cdot d\ps{1}\otimes d\ps{2}
\end{gather}
is bijective.
The map~\eqref{aux-cocan} is called the \emph{cocanonical map} of the coalgebra ${\cal B}$-coextension.

It is evident that the notion of \emph{algebra-Galois coalgebra coextension} dualizes the notion of
\emph{coalgebra-Galois algebra extension} by formal inverting all arrows in all diagrams and interchanging \emph{left}
and \emph{right}.
In the Hopf--Galois setting the construction dualizes the Hopf--Galois extensions~\cite{Schn93}.

\subsection{Formal Pontryagin duality and Takeuchi--Galois correspondence}\label{subsect-Pontryagin-dual}

In this subsection we will f\/irst summarize the basic properties of the dualization functor on vector spaces
from~\cite{BergmanHausknecht-book}.
We will then recall a~one-to-one correspondence between the coideal subalgebras and quotient coalgebras, known as the
Takeuchi--Galois correspondence~\cite{Take72}.

From the point of view of linear topological vector spaces with continuous linear mappings as morphisms, the dualization
functor def\/ines an equivalence between the opposite category of discrete vector spaces and the category of linearly
compact vector spaces~\cite[Proposition~24.8]{BergmanHausknecht-book}.
Moreover, it transforms naturally the algebraic tensor product of discrete vector spaces into a~completed one of
linearly compact spaces~\cite[Corollary~24.25]{BergmanHausknecht-book}.
In other words, dualization is a~strong monoidal functor.
Therefore one can regard on equal footing all structures def\/ined by diagrams in vector spaces together with their dual
counterparts obtained by reversing all arrows in all necessary diagrams.
This regards linear subspaces and quotient spaces, (co)algebras and Hopf algebras, as well as their bi(co)modules or one
sided (co)modules, one-sided and two-sided (co)ideals, and their (co)tensor products.
Hence, having a~diagrammatical proof of a~theorem in a~symmetric monoidal category of linear (topological) spaces, one
has automatically a~dual theorem after an appropriate dualization of the structures.
It is well known that the notion of Hopf algebra is self-dual.
In particular, it transforms the Hopf group-algebra of a~discrete group into a~linearly compact topological Hopf algebra
of functions on that group, customarily regarded as a~group algebra of a~dual compact quantum group.
Therefore we will call this duality \emph{Pontryagin} to separate it from cyclic duality.
We will show that the notion of SAYD module is Pontryagin self-dual up to the interchange of \emph{left} and
\emph{right}.
According to~\cite{Brze97,BrzeHaja99} the notion of coalgebra-Galois extension of algebras is Pontryagin dual to the
notion of algebra-Galois coextension of coalgebras.
Both dualities play a~role in the present paper.

For any Hopf algebra ${\cal H}$ with multiplication $\mu\colon {\cal H} \otimes{\cal H}\longrightarrow {\cal H}$, unit $\eta\colon
k\longrightarrow {\cal H}$, comultiplication $\Delta\colon {\cal H}\longrightarrow{\cal H}\otimes{\cal H}$, counit
$\varepsilon\colon {\cal H}\longrightarrow k$, and antipode $S\colon {\cal H}\longrightarrow{\cal H}$, Takeuchi introduces
in~\cite{Take72} the one-to-one (Galois) correspondence
\begin{gather*}
\{{\cal B}\subseteq {\cal H} \mid {\cal B}\ {\rm is\ a\ left\ coideal\ subalgebra}, {\cal H}\ {\rm is\ faithfully\ f\/lat\
over}\ {\cal B} \}
\\
\hspace*{50mm}\uparrow\downarrow
\\
\{{\cal I}\subseteq {\cal H} \mid {\cal I}\ {\rm is\ a\ coideal\ right\ ideal}, {\cal H}\ {\rm is\ faithfully\ cof\/lat\
over}\ {\cal H}/{\cal I} \},
\end{gather*}
under which ${\cal B}\mapsto {\cal B}^{+}{\cal H}$ and ${\cal I}\mapsto {\cal H}^{\co {\cal H}/{\cal I}}$.

A crucial observation for our purpose is that the above correspondence can be written categorically as an equivalence
between the category of left ${\cal H}$-comodule f\/lat extensions $i\colon {\cal B}\to{\cal H}$ and the category of right
${\cal H}$-module cof\/lat coextensions $\pi\colon {\cal H}\to {\cal C}$, which is given~by
\begin{gather*}%\label{diag}
{\cal B}\mapsto {\rm Coeq}({\cal H}\leftleftarrows {\cal B}\otimes{\cal H}),
\qquad
{\cal C}\mapsto {\rm Eq}({\cal H}\rightrightarrows {\cal H}\otimes{\cal C}),
\end{gather*}
where the parallel pair of left arrows in the coequalizer and the parallel pair of right arrows in the equalizer read as composites
\begin{gather*}
\xymatrix@R-1pc{{\cal H} & {\cal H}\otimes{\cal H}\ar[l]_-{\mu}& {\cal B}\otimes{\cal H}\ar[l]_-{\iota\otimes {\cal H}},
\\
{\cal H} & k\otimes {\cal H}\ar[l]_-{\cong} &{\cal B}\otimes{\cal H}\ar[l]_-{\varepsilon\iota\otimes {\cal H}},}\qquad
 \xymatrix@R-1pc{{\cal H}\ar[r]^-{\Delta} & {\cal H}\otimes{\cal H} \ar[r]^-{{\cal H}\otimes\pi} &{\cal H}\otimes{\cal C},
\\
{\cal H} \ar[r]^-{\cong} & {\cal H}\otimes k \ar[r]^-{{\cal H}\otimes \pi\eta} &{\cal H}\otimes{\cal C},}
\end{gather*}
respectively.
It is thus evident that this correspondence is Pontryagin self-dual up to an interchange of \emph{left} and \emph{right}
in dual structures.
We will say that ${\cal B}$ and ${\cal C}$ as above are \emph{Takeuchi--Galois transforms} of each other.

\subsection{Relative cyclic homology of (co)algebra (co)extensions}

We start with a~quick detour on the relative cyclic homology of algebra extensions.
For any ${\cal B}$-bimodule ${\cal M}$ we let
\begin{gather*}
[{\cal M}]_{\cal B}:={\cal M}/[{\cal M},{\cal B}],
\end{gather*}
where $[{\cal M},{\cal B}]$ is the subspace of ${\cal M}$ generated by all commutators $[m,b]:=m\cdot b-b\cdot m$.
We will denote by $[m]_{\cal B}$ the class of $m\in {\cal M}$ in $[{\cal M}]_{\cal B}$.

The relative cyclic homology of an algebra extension ${\cal B}\rightarrow {\cal A}$ is computed by the cyclic object
\begin{gather}
\label{aux-rel-cyclic-alg}
{\rm C}_n({\cal A}\mid{\cal B}):=[{\cal A}^{\otimes_{\cal B} n+1}]_{\cal B},
\end{gather}
equipped, for all $n\geq 0$, with the morphisms
\begin{gather*}%\label{aux-faces-Z-n(A/B, A)}
d_i\colon \ {\rm C}_n({\cal A}\mid{\cal B}) \longrightarrow {\rm C}_{n-1}({\cal A}\mid{\cal B}),
\qquad
0 \leq j \leq n,
\\
d_i\big[a^0 \otimes_{\cal B} a^1 \otimes_{\cal B} \dots \otimes_{\cal B} a^n\big]_{\cal B}\\
\qquad{}
=\begin{cases}
\big[a^0 a^1 \otimes_{\cal B} \dots \otimes_{\cal B} a^n \big]_{\cal B}, & i =0,
\\
\big[a^0 \otimes_{\cal B} a^1 \otimes_{\cal B} \dots \otimes_{\cal B} a^i a^{i+1} \otimes_{\cal B} \dots \otimes_{\cal B} a^n \big]_{\cal B},& 1\leq i\leq n-1,
\\
\big[a^n a^0 \otimes_{\cal B} a^1 \otimes_{\cal B} \dots \otimes_{\cal B} a^{n-1}\big]_{\cal B}, & i =n,
\end{cases}
\\
%\label{aux-degeneracies-Z-n(A/B, A)}
s_i\colon \ {\rm C}_n({\cal A}\mid{\cal B}) \longrightarrow {\rm C}_{n+1}({\cal A}\mid{\cal B}),
\qquad
0 \leq j \leq n,
\\
s_i\big[a^0 \otimes_{\cal B} \dots \otimes_{\cal B} a^n\big]_{\cal B} =\big[a^0 \otimes_{\cal B} a^1
\otimes_{\cal B} \dots \otimes_{\cal B} a^i \otimes_{\cal B} 1 \otimes_{\cal B} a^{i+1} \otimes_{\cal B} \dots
\otimes_{\cal B} a^n \big]_{\cal B}
\end{gather*}
and
\begin{gather*}%\label{aux-cyclic-Z-n(A/B, A)}
t_n\colon \ {\rm C}_n({\cal A}\mid{\cal B}) \longrightarrow {\rm C}_n({\cal A}\mid{\cal B}),
\qquad
t_n\big[a^0 \otimes_{\cal B} \dots \otimes_{\cal B} a^n\big]_{\cal B} =\big[a^n \otimes_{\cal B} a^0
\otimes_{\cal B} \dots \otimes_{\cal B} a^{n-1}\big]_{\cal B}.
\end{gather*}
Hochschild homology of the complex~\eqref{aux-rel-cyclic-alg} is called the Hochschild homology of the extension, and is
denoted by ${\rm HH}_\bullet({\cal A}\mid{\cal B})$.
Similarly, cyclic (resp.\
periodic cyclic, negative cyclic) homology of the cyclic object~\eqref{aux-rel-cyclic-alg} is called the relative cyclic
(resp.\ Hochschild, periodic cyclic, negative cyclic) homology of the extension ${\cal B} \rightarrow{\cal A}$, and it is
denoted by ${\rm HC}_\bullet({\cal A}\mid{\cal B})$ (resp.\
${\rm HH}_\bullet({\cal A}\mid{\cal B})$, ${\rm HP}_\bullet({\cal A}\mid{\cal B})$, ${\rm HN}_\bullet({\cal A}\mid{\cal B})$).

The natural mapping
\begin{gather*}
{\rm C}_n({\cal A})\longrightarrow {\rm C}_n({\cal A}\mid{\cal B}),
\qquad
a^0\otimes\dots\otimes a^n\mapsto \big[a^0\otimes_{\cal B} \dots \otimes_{\cal B} a^n\big]_{{\cal B}},
\qquad
n\geq 0,
\end{gather*}
induces a~map ${\rm HH}_\bullet({\cal A})\longrightarrow {\rm HH}_\bullet({\cal A}\mid{\cal B})$ of Hochschild homology
groups, which further induces a~mapping ${\rm HC}_\bullet({\cal A})\longrightarrow {\rm HC}_\bullet({\cal A}\mid{\cal
B})$ of cyclic homology groups.
This is the computational motivation behind the relative cyclic homology~\cite{Kadi92}, in which it is proved to be an
isomorphism when ${\cal B}\subseteq {\cal A}$ is a~separable subalgebra (semisimple in characteristic zero).
The relative cyclic homology then appeared in~\cite{Scha92}, where the relative cyclic homology ${\rm HC}_\bullet(kG\mid kN)$
of group algebras, associated to normal subgroups, was computed, and is applied to extend Eckmann's
result~\cite{Eckm86} on the Bass conjecture.

Dually, for any ${\cal C}$-bicomodule ${\cal M}$ one def\/ines ${\cal M}^{\cal C}:=\{m\in {\cal M}\mid m\ns{0}\otimes
m\ns{1} =m\ns{0}\otimes m\ns{-1}\}$. Then the relative cyclic cohomology of a~coalgebra coextension ${\cal
D}\rightarrow {\cal C}$ is computed by the cocyclic object
\begin{gather}
\label{aux-rel-cyclic-coalg}
{\rm C}^n({\cal D}\mid{\cal C}):=\big({\cal D}^{\Box_{\cal C} n+1}\big)^{\cal C},
\qquad
n\geq 0,
\end{gather}
with the structure maps
\begin{gather*}%\label{aux-cofaces-Z-n(C/D, C)}
\delta_i\colon \ {\rm C}^n({\cal D}\mid{\cal C}) \longrightarrow {\rm C}^{n+1}({\cal D}\mid{\cal C}),
\qquad
0 \leq j \leq n+1,
\\
\delta_i\big(d^0\otimes d^1\otimes \cdots\otimes d^n\big)
=\begin{cases}
d^0\ps{1}\otimes d^0\ps{2}\otimes d^1\otimes \cdots\otimes d^n, & i =0,
\\
d^0\otimes \cdots \otimes  d^i\ps{1}\otimes d^i\ps{2}\otimes \cdots \otimes d^n, & 1 \leq i\leq n,
\\
d^0\ps{2}\otimes d^1\otimes \cdots\otimes d^n\otimes d^0\ps{1}, & i =n+1,
\end{cases}
\\
%\label{aux-codegeneracies-Z-n(C/D, C)}
\sigma_j\colon \ {\rm C}^n({\cal D}\mid{\cal C}) \longrightarrow {\rm C}^{n-1}({\cal D}\mid{\cal C}),
\qquad
0 \leq j \leq n-1,
\\
\sigma_j\big(d^0\otimes d^1\otimes \cdots\otimes d^n\big) =d^0\otimes d^1\otimes \cdots \otimes
d^{j}\varepsilon\big(d^{j+1}\big)\otimes \cdots \otimes d^n
\end{gather*}
and
\begin{gather*}%\label{aux-cocyclic-Z-n(C/D, C)}
\tau_n\colon \ {\rm C}^n({\cal D}\mid{\cal C}) \longrightarrow {\rm C}^n({\cal D}\mid{\cal C}),
\qquad
\tau_n\big(d^0\otimes d^1\otimes \cdots\otimes d^n\big) =d^1\otimes \cdots\otimes d^n\otimes d^0.
\end{gather*}
Then the cyclic duality~\cite{Conn83,KhalRang05} yields the cyclic module structure on the collection of vector spaces
\begin{gather}
\label{cyc-mod-coext}
{\rm C}_n({\cal D}\mid{\cal C}):=\big({\cal D}^{\Box_{\cal C} n+1}\big)^{\cal C}
\end{gather}
given by faces
\begin{gather*}%\label{aux-cofaces-Z-n(C/D, C)-cyclic-dual}
d_i\colon \ {\rm C}_n({\cal D}\mid{\cal C}) \longrightarrow {\rm C}_{n-1}({\cal D}\mid{\cal C}),
\qquad
0 \leq i \leq n,
\\
d_i\big(d^0\otimes d^1\otimes \cdots\otimes d^n\big) =d^0\otimes d^1\otimes \cdots \otimes \varepsilon\big(d^i\big)d^{i+1}\otimes\cdots \otimes d^n,
\qquad
0 \leq i \leq n-1,
\\
d_n\big(d^0\otimes d^1\otimes \cdots\otimes d^n\big) =d^0\otimes d^1\otimes \cdots \otimes d^{n-1}\varepsilon\big(d^n\big),
\end{gather*}
degeneracies
\begin{gather*}%\label{aux-codegeneracies-Z-n(C/D, C)-cyclic-dual}
s_i\colon \ {\rm C}_n({\cal D}\mid{\cal C}) \longrightarrow {\rm C}_{n+1}({\cal D}\mid{\cal C}),
\qquad
0 \leq i \leq n,
\\
s_i\big(d^0\otimes d^1\otimes \cdots\otimes d^n\big) =d^0\otimes d^1\otimes \cdots \otimes\Delta\big(d^i\big)\otimes \cdots \otimes
d^n,
\end{gather*}
and the cyclic operator
\begin{gather*}%\label{aux-cocyclic-Z-n(C/D, C)-cyclic-dual}
t_n\colon \ {\rm C}_n({\cal D}\mid{\cal C}) \longrightarrow {\rm C}_n({\cal D}\mid{\cal C}),
\qquad
t_n\big(d^0\otimes d^1\otimes \cdots\otimes d^n\big) =d^n\otimes d^0\otimes \cdots\otimes d^{n-1}.
\end{gather*}
Hochschild homology of the complex~\eqref{cyc-mod-coext} is called the Hochschild homology of the coextension, and is
denoted by ${\rm HH}_\bullet({\cal D}\mid{\cal C})$.
Similarly, cyclic (resp.\
periodic cyclic, negative cyclic) homology of the cyclic object~\eqref{aux-rel-cyclic-coalg} is called the relative
cyclic (resp.\
Hochschild, periodic cyclic, negative cyclic) homology of the coextension ${\cal D} \rightarrow{\cal C}$, and it is
denoted by ${\rm HC}_\bullet({\cal D}\mid{\cal C})$ (resp.\
${\rm HH}_\bullet({\cal D}\mid{\cal C})$, ${\rm HP}_\bullet({\cal D}\mid{\cal C})$, ${\rm HN}_\bullet({\cal D}\mid{\cal C})$).

\subsection[Hopf-cyclic homology of ${\cal H}$-module coalgebras and ${\cal H}$-comodule algebras]{Hopf-cyclic homology
of $\boldsymbol{\cal H}$-module coalgebras and $\boldsymbol{\cal H}$-comodule algebras}

In this subsection we recall the relative Hopf-cyclic homology with coef\/f\/icients, for coalgebras, using the cyclic
duality principle~\cite{Conn83,KhalRang05}.

Let us f\/irst recall the def\/inition of a~left-right stable anti-Yetter--Drinfeld module over a~Hopf algebra ${\cal H}$
from~\cite{HajaKhalRangSomm04-I}.
Let a~linear space ${\cal M}$ be a~left ${\cal H}$-module by $\ell\colon {\cal H}\otimes{\cal M}\longrightarrow {\cal M}$
given by $\ell(h\otimes m)=h\cdot m$, and a~right ${\cal H}$-comodule via $\rho\colon {\cal M} \longrightarrow {\cal M}
\otimes {\cal H}$ given by $\rho(m)=m\ns{0} \otimes m\ns{1}$.
Then ${\cal M}$ is called a~left-right anti-Yetter--Drinfeld module (AYD module) over ${\cal H}$ if the ${\cal
H}$-action and ${\cal H}$-coaction are compatible as
\begin{gather*}
(h \cdot m)\ns{0}\otimes (h \cdot m)\ns{1} =h\ps{2} \cdot m\ns{0} \otimes h\ps{3} m\ns{1} S(h\ps{1}).
\end{gather*}
A~left-right AYD module ${\cal M}$ is called \emph{stable} (and then is abbreviated as SAYD module) if
\begin{gather*}
m\ns{1}\cdot m\ns{0} =m.
\end{gather*}
On the other hand, a~right-left SAYD module structure is def\/ined in terms of a~right ${\cal H}$-module structure via
$r\colon {\cal M}\otimes{\cal H}\longrightarrow {\cal M}$ given by $r(m\otimes h)=m\cdot h$, and a~left ${\cal H}$-comodule
structure via $\lambda\colon {\cal M} \longrightarrow {\cal H} \otimes {\cal M}$ given by $\lambda(m)=m\ns{-1} \otimes
m\ns{0}$ such that
\begin{gather*}
(m \cdot h)\ns{-1}\otimes (m \cdot h)\ns{0} =S(h\ps{3}) \cdot m\ns{-1} h\ps{1} \otimes m\ns{0} h\ps{2},
\qquad
m\ns{0}\cdot m\ns{-1} =m.
\end{gather*}
Note that these conditions are Pontryagin dual to each other.
Indeed, the left-right SAYD module compatibility reads as commutativity of the following diagrams in the category $V\
ect$ of vector spaces ($\sigma$ being the transposition of the corresponding tensorands)
\begin{gather*}
\xymatrix{{\cal H}\otimes {\cal H}\otimes {\cal H}\otimes {\cal M}\otimes {\cal H}\ar[rr] && {\cal H}\otimes {\cal
M}\otimes {\cal H}\otimes {\cal H}\otimes {\cal H}\ar[d]^-{\ell \otimes \mu_{3}}
\\
{\cal H}\otimes {\cal M}\ar[u]^-{\Delta^3\otimes \rho}\ar[r]^-{\ell} & {\cal M} \ar[r]^-{\rho}&{\cal M}\otimes {\cal H},}
\end{gather*}
where the upper horizontal arrow admits two decompositions
\begin{gather*}
\xymatrix{& {\cal H}\otimes {\cal H}\otimes {\cal M}\otimes {\cal H}\otimes {\cal H} \ar[rd]^{\  \ \
H\otimes\sigma_{{\cal H},{\cal M}}\otimes {\cal H}\otimes S}&
\\
{\cal H}\otimes {\cal H}\otimes {\cal H}\otimes {\cal M}\otimes {\cal H}\ar[ur]^{\sigma_{{\cal H}, {\cal H}\otimes {\cal
H}\otimes {\cal M}\otimes {\cal H}}\ \ \ \ }\ar[dr]_{S\otimes {\cal H} \otimes \sigma_{{\cal H}, {\cal
M}}\otimes {\cal H} \ \ \ \ \ } && {\cal H}\otimes {\cal M}\otimes {\cal H}\otimes {\cal H}\otimes {\cal H}
\\
&{\cal H}\otimes {\cal H}\otimes {\cal M}\otimes {\cal H}\otimes {\cal H}\ar[ru]_{\ \ \ \ \ \sigma_{{\cal H}\otimes
{\cal H}\otimes {\cal M}\otimes {\cal H}, {\cal H}}},}
\end{gather*}
and
\begin{gather*}
\xymatrix{{\cal M}\otimes {\cal H} \ar[r]^-{\sigma_{{\cal M},{\cal H}}} & {\cal H}\otimes {\cal M}\ar[d]^-{\ell}
\\
{\cal M}\ar[u]^{\rho}\ar[r]^{=} & {\cal M}.}
\end{gather*}
Now we see that reversing the arrows, inverting transpositions~$\sigma$ and next interchanging in pairs~$\Delta$
and~$\mu$, $\ell $ and~$\lambda$,~$\rho$ and~$r$, and f\/inally \emph{left} and \emph{right}, we obtain the right-left
SAYD module compatibility.

Let us now recall the Hopf-cyclic cohomology of a~Hopf-module coalgebra with coef\/f\/icients.
Let ${\cal H}$ be a~Hopf algebra with an invertible antipode, ${\cal C}$ a~right ${\cal H}$-module coalgebra, i.e.,
a~right ${\cal H}$-module such that
\begin{gather*}
\Delta(c\cdot h) =c\ps{1}\cdot h\ps{1}\otimes c\ps{2}\cdot h\ps{2},
\qquad
\varepsilon(c\cdot h) =\varepsilon(c)\varepsilon(h),
\qquad
\forall\,  c\in {\cal C}, \quad \forall\,  h\in{\cal H},
\end{gather*}
and ${\cal M}$ a~left-right SAYD module over ${\cal H}$.
Then for all $n\geq 0$ we def\/ine
\begin{gather}
\label{aux-cyclic-module-C-H-M}
{\rm C}^n({\cal C},{\cal M})_{{\cal H}}:={\cal C}^{\otimes n+1} \otimes_{{\cal H}}{\cal M},
\end{gather}
where on the right hand side the diagonal right ${\cal H}$-module structure on ${\cal C}^{\otimes n+1}$ is used.

{\sloppy The subscript ${\cal H}$ refers to the coinvariants of the diagonal right ${\cal H}$-module structure \mbox{${\cal C}^{\otimes
n+1}\otimes{\cal M}$}, by switching the one on ${\cal M}$ from left to right via the antipode.
The collection of vector spaces ${\rm C}^n({\cal C},{\cal M})_{{\cal H}}$ together with the operators
\begin{gather*}
\delta_i\colon \ {\rm C}^{n-1}({\cal C},{\cal M})_{{\cal H}} \to {\rm C}^n({\cal C},{\cal M})_{{\cal H}},
\qquad
0 \leq i \leq n,
\\
\delta_i\big(\big(c^0 \otimes \cdots \otimes c^{n-1}\big) \otimes_{{\cal H}} m\big) \\
\qquad{} =\big(c^0 \otimes \cdots \otimes c^i\ps{1} \otimes
c^i\ps{2} \otimes \cdots \otimes c^{n-1}\big) \otimes_{{\cal H}} m,
\qquad
0 \leq i \leq n-1,
\\
\delta_n\big(\big(c^0 \otimes \cdots \otimes c^{n-1}\big) \otimes_{{\cal H}} m\big) =\big(c^0\ps{2} \otimes \cdots \otimes c^{n-1}
\otimes c^0\ps{1} \cdot S^{-1}(m\ns{1})\big) \otimes_{{\cal H}} m\ns{0},
\\
\sigma_j\colon \ {\rm C}^{n+1}({\cal C},{\cal M})_{{\cal H}} \to {\rm C}^n({\cal C},{\cal M})_{{\cal H}},
\qquad
0 \leq i \leq n,
\\
\sigma_j\big(\big(c^0 \otimes \cdots \otimes c^{n+1}\big) \otimes_{{\cal H}} m\big) =\big(c^0 \otimes \cdots \otimes
c^{j}\varepsilon\big(c^{j+1}\big) \otimes \cdots \otimes c^{n+1}\big) \otimes_{{\cal H}} m,
\end{gather*}
and
\begin{gather*}
\tau_n\colon \ {\rm C}^{n}({\cal C},{\cal M})_{{\cal H}} \to {\rm C}^n({\cal C},{\cal M})_{{\cal H}},
\\
\tau_n\big(\big(c^0 \otimes \cdots \otimes c^n\big) \otimes_{\mathcal{H}} m\big) =\big(c^1 \otimes \cdots \otimes c^n \otimes c^0 \cdot
S^{-1}(m\ns{1})) \otimes_{\mathcal{H}} m\ns{0}.
\end{gather*}
is a~cocyclic module.
Cyclic cohomology of this cocyclic module is called the Hopf-cyclic cohomology of the ${\cal H}$-module coalgebra ${\cal
C}$ with coef\/f\/icients in the left-right SAYD module ${\cal M}$ over~${\cal H}$, and is denoted by ${\rm
HC}^\bullet({\cal C},{\cal M})_{\cal H}$.

}

Applying the cyclic duality procedure~\cite{Conn83,KhalRang05} we obtain on~\eqref{aux-cyclic-module-C-H-M} the cyclic
module structure given by the faces
\begin{gather*}%\label{cyclic-X-face}
d_i\colon \ {\rm C}_{n+1}({\cal C},{\cal M})_{{\cal H}} \to {\rm C}_n({\cal C},{\cal M})_{{\cal H}},
\qquad
0 \leq i \leq n+1,
\\
d_i\big(\big(c^0 \otimes \cdots \otimes c^{n+1}) \otimes_{{\cal H}} m\big) =\big(c^0 \otimes \cdots \otimes \varepsilon\big(c^i\big)c^{i+1}
\otimes \cdots \otimes c^{n+1}\big) \otimes_{{\cal H}} m,
\qquad
0 \leq i \leq n,
\\
d_{n+1}\big(\big(c^0 \otimes \cdots \otimes c^{n+1}\big) \otimes_{{\cal H}} m\big) =\big(\varepsilon\big(c^{n+1}\big) c^0 \otimes \cdots \otimes
c^{n}\big) \otimes_{{\cal H}} m
\end{gather*}
degeneracies
\begin{gather*}%\label{cyclic-X-degeneracy}
s_i\colon \ {\rm C}_{n-1}({\cal C},{\cal M})_{{\cal H}} \to {\rm C}_n({\cal C},{\cal M})_{{\cal H}},
\qquad
0 \leq i \leq n-1,
\\
s_i\big(\big(c^0 \otimes \cdots \otimes c^{n-1}\big) \otimes_{\mathcal{H}} m\big) =\big(c^0 \otimes \cdots \otimes c^i\ps{1} \otimes
c^i\ps{2} \otimes \cdots \otimes c^{n-1}\big) \otimes_{\mathcal{H}} m,
\end{gather*}
and the cyclic operator
\begin{gather*}%\label{cyclic-X-cyclic}
t_n\colon \ {\rm C}_{n}({\cal C},{\cal M})_{{\cal H}} \to {\rm C}_n({\cal C},{\cal M})_{{\cal H}},
\\
t_n\big(\big(c^0 \otimes \cdots \otimes c^n\big) \otimes_{{\cal H}} m\big) =\big(c^n \cdot m\ns{1} \otimes c^0 \otimes \cdots \otimes
c^{n-1}\big) \otimes_{{\cal H}} m\ns{0}.
\end{gather*}
The cyclic homology of this cyclic module is called the cyclic (resp.\
Hochschild, periodic cyclic, negative cyclic) homology of the $\mathcal{H}$-module coalgebra ${\cal C}$, with
coef\/f\/icients in the left-right SAYD module ${\cal M}$ over ${\cal H}$, and is denoted by ${\rm HC}_\bullet({\cal
C},{\cal M})_{\cal H}$ (resp.\
${\rm HH}_\bullet({\cal C},{\cal M})_{\cal H}$, ${\rm HP}_\bullet({\cal C},{\cal M})_{\cal H}$, ${\rm HN}_\bullet({\cal
C},{\cal M})_{\cal H}$).

We conclude this section with the Hopf-cyclic homology of ${\cal H}$-comodule algebras~\cite{HajaKhalRangSomm04-II}.
Let ${\cal M}$ be a~left-right SAYD module over the Hopf algebra ${\cal H}$, and ${\cal B}$ a~left ${\cal H}$-comodule algebra.
Then for $n\geq 0$ we def\/ine ${\rm C}_{n}({\cal B},{\cal M})^{{\cal H}}:={\cal M} \Box_{\cal H} {\cal B}^{\otimes n+1}$,
where on the right hand side the diagonal left ${\cal H}$-comodule structure on ${\cal B}^{\otimes n+1}$ is used.
The superscript ${\cal H}$ refers to the invariants of the diagonal right ${\cal H}$-comodule structure ${\cal
M}\otimes{\cal B}^{\otimes n+1}$, switching the one on ${\cal B}^{\otimes n+1}$ from left to right by means of the
antipode.

The collection of these vector spaces together with operators
\begin{gather*}
d_i\colon \ {\rm C}_n({\cal B},{\cal M})^{{\cal H}} \longrightarrow {\rm C}_{n-1}({\cal B},{\cal M})^{{\cal H}},
\qquad
0 \leq i \leq n,
\\
d_i\big(m\otimes b^0 \otimes \cdots \otimes b^n\big) =m\otimes b^0 \otimes \cdots \otimes b^ib^{i+1}\otimes \cdots \otimes b^n,
\qquad
0 \leq i \leq n-1,
\\
d_n\big(m\otimes b^0 \otimes \cdots \otimes b^n\big) =b^n\ns{-1}\cdot m\otimes b^n\ns{0}b^0\otimes b^1 \otimes \cdots \otimes b^{n-1},
\\
s_j\colon \ {\rm C}_n({\cal B},{\cal M})^{{\cal H}} \longrightarrow {\rm C}_{n+1}({\cal B},{\cal M})^{{\cal H}},
\qquad
0 \leq i \leq n,
\\
s_j\big(m\otimes b^0 \otimes \cdots \otimes b^n\big) =m\otimes b^0 \otimes\cdots \otimes b^j\otimes 1\otimes\cdots \otimes b^n,
\end{gather*}
and
\begin{gather*}
t_n\colon \ {\rm C}_n({\cal B},{\cal M})^{{\cal H}} \longrightarrow {\rm C}_n({\cal B},{\cal M})^{{\cal H}},
\\
t_n\big(m\otimes b^0 \otimes \cdots \otimes b^n\big) =b^n\ns{-1}\cdot m\otimes b^n\ns{0}\otimes b^0\otimes \cdots \otimes b^{n-1}
\end{gather*}
is a~cyclic module.
Finally, the cyclic homology of this cyclic module is called the cyclic (resp.\
Hochschild, periodic cyclic, negative cyclic) homology of the $\mathcal{H}$-comodule algebra ${\cal C}$, with
coef\/f\/icients in the left-right SAYD module ${\cal M}$ over ${\cal H}$, and is denoted by ${\rm HC}_\bullet({\cal
C},{\cal M})^{\cal H}$ (resp.\
${\rm HH}_\bullet({\cal C},{\cal M})^{\cal H}$, ${\rm HP}_\bullet({\cal C},{\cal M})^{\cal H}$, ${\rm HN}_\bullet({\cal
C},{\cal M})^{\cal H}$).

\section{Relative cyclic homology as Hopf-cyclic homology\\ with coef\/f\/icients}

In this section we achieve our main result identifying the relative cyclic homology of a~homogeneous ${\cal C}$-Galois
extension ${\cal B}\subseteq{\cal H}$ with the Hopf-cyclic homology, with coef\/f\/icients, of the right ${\cal H}$-module
coalgebra~${\cal C}$.
Moreover, in view of the Pontryagin duality of Subsection~\ref{subsect-Pontryagin-dual}, we obtain an identif\/ication of
the relative cyclic homology of a~${\cal B}$-Galois coextension ${\cal H}\to {\cal C}$ with the Hopf-cyclic homology,
with (the Pontryagin dual) coef\/f\/icients, of the ${\cal H}$-comodule algebra ${\cal B}$.
We shall conclude the section developing spectral sequences to shed further light on the relative homology groups.

\subsection{The isomorphism of cyclic objects}\label{subsect-iso-cyclic-obj}

In this subsection we will construct an explicit isomorphism from the relative homology complex of a~homogeneous ${\cal
H}/{\cal I}$-Galois extension ${\cal B}:={\cal H}^{\co {\cal H}/{\cal I}}\to {\cal H}$ to the Hopf-cyclic homology
complex of the ${\cal H}$-module coalgebra ${\cal H}/{\cal I}$.

Let ${\cal H}$ be a~Hopf algebra and ${\cal I} \subseteq {\cal H}$ a~coideal right ideal of ${\cal H}$.
Then ${\cal H}/{\cal I}$ becomes a~coaugmented quotient coalgebra in a~canonical way,
\begin{gather*}%\label{comul}
\Delta(\wbar{h}):=\overline{h\ps{1}}\otimes \overline{h\ps{2}},
\qquad
\varepsilon(\wbar{h}):=\varepsilon(h),
\qquad
\forall\,  h\in{\cal H},
\end{gather*}
where $\wbar{h}:=h+{\cal I}$, and the canonical coaugmentation of ${\cal H}/{\cal I}$ is given by the group-like
$\overline{1}$.

Let ${\cal B} \subseteq {\cal H}$ be a~homogeneous ${\cal H}/{\cal I}$-Galois extension given by the canonical right
${\cal H}/{\cal I}$-coaction
\begin{gather*}%\label{aux-homogeneous H/I-extension}
\xymatrix{{\cal H}\ar[r]^{\Delta} & {\cal H}\otimes {\cal H}\ar[r] &{\cal H} \otimes ({\cal H}/{\cal I}),
\qquad
h \mapsto h\ps{1} \otimes \wbar{h\ps{2}}}
\end{gather*}
on ${\cal H}$.
In view of the def\/inition ${\cal B}:={\cal H}^{\co {\cal H}/{\cal I}}$, which reads as
\begin{gather}
\label{coinv}
b\ps{1}\otimes \overline{b\ps{2}} =b\otimes \overline{1},
\end{gather}
applying the counit to the left tensorands we deduce $\overline{b}=\varepsilon(b)\overline{1}$.
We will use also the following iterated form of~\eqref{coinv}:
\begin{gather}
\label{coinviter}
b\ps{1}\otimes \dots\otimes\overline{b\ps{n+2}} =b\ps{1}\otimes \dots\otimes b\ps{n+1}\otimes\overline{1}.
\end{gather}

The three following facts are crucial to our purposes.
First of all, ${\cal B}$ is the left coideal of~${\cal H}$~\cite{BrzeHaja09}, i.e.,
\begin{gather*}%\label{coid-sub}
\Delta({\cal B})\subset {\cal H}\otimes {\cal B}.
\end{gather*}
Secondly, the Galois condition f\/ixes the coideal right ideal ${\cal I}\subseteq {\cal H}$ completely as
follows~\cite[Theorem~2.6]{BrzeHaja09}.

\begin{Theorem}
%\label{thm-I=B+H}
Let ${\cal B}\subseteq {\cal H}$ be a~homogeneous ${\cal H}/{\cal I}$-extension.
Then this extension is Galois if and only if ${\cal I} ={\cal B}^{+}{\cal H}$, where ${\cal B}^{+}:={\cal B} \cap \operatorname{Ker}
\varepsilon$.
\end{Theorem}

Finally, there is an explicit formula for the translation map, and hence for the inverse to the canonical map,~\cite[Corollary~2.8]{BrzeHaja09}.

\begin{Corollary}
Let ${\cal B}\subseteq {\cal H}$ be a~coalgebra-Galois ${\cal H}/{\cal I}$-extension as above.
Then the translation map $\tau:=\operatorname{can}^{-1}(1 \otimes -)$ is given~by
\begin{gather}
\label{trans}
\tau(\overline{h}) =S(h\ps{1}) \otimes_{B} h\ps{2}.
\end{gather}
\end{Corollary}

Moreover, it implies that the right hand side of~\eqref{trans} is independent of the choice of the representative~$h$ of
the class $\overline{h}$.

The canonical map
\begin{gather*}%\label{can}
{\cal H}\otimes_{{\cal B}}{\cal H}\rightarrow {\cal H}\otimes{\cal H}/{\cal I},
\qquad
h\otimes_{{\cal B}}h' \mapsto hh'\ps{1}\otimes \overline{h'\ps{2}}
\end{gather*}
and its inverse
\begin{gather*}%\label{can-inverse}
{\cal H}\otimes{\cal H}/{\cal I} \rightarrow {\cal H}\otimes_{{\cal B}}{\cal H},
\qquad
h\otimes \overline{h'}\mapsto h S(h'\ps{1})\otimes_{{\cal B}} h'\ps{2}
\end{gather*}
can be inductively extended to ${\cal H}$-bimodule isomorphisms
\begin{gather*}%\label{aux-can-n-map-H/I}
\operatorname{can}_n\colon \ {\cal M}\otimes_{\cal B}\otimes{\cal H}^{\otimes_{\cal B} n} \longrightarrow {\cal M} \otimes {\cal C}^{\otimes n},
\\
\operatorname{can}_n\big(m \otimes_{\cal B} h^1\otimes_{\cal B} \cdots \otimes_{\cal B} h^n\big)
\\
\qquad{}
=mh^1\ps{1}\dots h^n\ps{1} \otimes \wbar{h^1\ps{2}\cdots h^n\ps{2}} \otimes\cdots\otimes \wbar{h^{n-1}\ps{n}h^n\ps{n}}
\otimes \wbar{h^n\ps{n+1}}.
\end{gather*}
and
\begin{gather*}%\label{aux-inverse-can-n-map-H/I}
\operatorname{can}^{-1}_n\colon \  {\cal M} \otimes {\cal C}^{\otimes n} \longrightarrow {\cal M}\otimes_{\cal B} {\cal H}^{\otimes_{\cal B} n},
\\
\operatorname{can}^{-1}_n\big(m \otimes \wbar{g^1} \otimes\cdots\otimes \wbar{g^n}\big)
\\
\qquad{}
=mS\big(g^1\ps{1}\big) \otimes_{\cal B} g^1\ps{2}S\big(g^2\ps{1}\big) \otimes_{\cal B} \cdots \otimes_{\cal B} g^{n-1}\ps{2}S\big(g^n\ps{1}\big)
\otimes_{\cal B} g^n\ps{2},
\end{gather*}
respectively, for any ${\cal H}$-bimodule ${\cal M}$, see also~\cite[Proposition~3.6]{Kadi14}.

We also recall that ${\cal H}$, equipped with the (left) adjoint ${\cal H}$-action
\begin{gather}
\label{HonM}
h\triangleright h':=h\ps{2}h'S(h\ps{1})
\end{gather}
and the (right) ${\cal H}$-coaction given by the comultiplication, is a~left-right SAYD module over ${\cal H}$, which is
denoted by $\ad({\cal H})$, see for instance~\cite[Example~4.3]{JaraStef06}.
Moreover, this action satisf\/ies
\begin{gather}
\label{HonH}
h\ps{2}\triangleright (h'h\ps{1}) =h\ps{2}\ps{2}h'h\ps{1}S(h\ps{2}\ps{1})
=h\ps{3}h'h\ps{1}S(h\ps{2}) =h\ps{2}h'\varepsilon(h\ps{1}) =hh'.
\end{gather}

All that is used in the following lemma.

\begin{Lemma}
Let ${\cal B}\subseteq {\cal H}$ be a~homogeneous ${\cal H}/{\cal I}$-Galois extension.
Then for any $n \geq 0$ we have the isomorphism of vector spaces $[{\cal H}^{\otimes_{\cal B} n+1}]_{{\cal
B}} \cong ({\cal H}/{\cal I})^{\otimes n+1} \otimes_{\cal H} \ad({\cal H})$ implemented~by
\begin{gather}
\psi_n\colon \  ({\cal H}/{\cal I})^{\otimes n+1} \otimes_{\cal H} \ad({\cal H})\stackrel{\cong}{\longrightarrow}
\big[{\cal H}^{\otimes_{\cal B} n+1}\big]_{\cal B},
\label{aux-map-psi}
\\
\psi_n\big(\big(\wbar{g^0}\otimes\cdots\otimes \wbar{g^n}\big)\otimes_{\cal H} h \big)
=\big[g^n\ps{2}hS\big(g^0\ps{1}\big) \otimes_{\cal B} g^0\ps{2}S\big(g^1\ps{1}\big) \otimes_{\cal B} \cdots \otimes_{\cal B}g^{n-1}\ps{2}S\big(g^n\ps{1}\big)\big]_{\cal B},
\nonumber
\end{gather}
with the inverse
\begin{gather}
\varphi_n\colon \ \big[{\cal H}^{\otimes_{\cal B} n+1}\big]_{{\cal B}} \stackrel{\cong}{\longrightarrow} ({\cal H}/{\cal
I})^{\otimes n+1} \otimes_{\cal H} \ad({\cal H}),
\nonumber
\\
\varphi_n\big(\big[h^0\otimes_{\cal B}\cdots\otimes_{\cal B} h^n\big]_{\cal B} \big)\nonumber\\
\qquad{}
=\big(\wbar{h^1\ps{2}\cdots h^n\ps{2}} \otimes\cdots\otimes \wbar{h^{n-1}\ps{n}h^n\ps{n}}
\otimes
\wbar{h^n\ps{n+1}} \otimes \wbar{1}\big) \otimes_{\cal H} h^0h^1\ps{1}\cdots h^n\ps{1}.
\label{aux-map-vp}
\end{gather}
\end{Lemma}

\begin{proof}
First of all, we have to prove that the maps are well def\/ined.
Let us begin with~\eqref{aux-map-vp}.
We observe that
\begin{gather*}
\varphi_n\big(\big[h^0\otimes_{\cal B}\cdots\otimes_{\cal B} h^nb\big]_{\cal B}\big)
\\
\qquad
=\big(\wbar{h^1\ps{2}\cdots (h^nb)\ps{2}}
\otimes\dots\otimes \wbar{h^{n-1}\ps{n}(h^nb)\ps{n}} \otimes
\wbar{(h^nb)\ps{n+1}} \otimes \wbar{1} \big)
\otimes_{\cal H} h^0h^1\ps{1}\cdots (h^nb)\ps{1}
\\
\qquad{}
=\big(\wbar{h^1\ps{2}\cdots h^n\ps{2}b\ps{2}}
\otimes\dots\otimes \wbar{h^{n-1}\ps{n}h^n\ps{n}b\ps{n}} \otimes
\wbar{h^n\ps{n+1}b\ps{n+1}} \otimes \wbar{1} \big)\!
\otimes_{\cal H} h^0h^1\ps{1}\cdots h^n\ps{1}b\ps{1}
\\
\qquad
=\big(\wbar{h^1\ps{2}\cdots h^n\ps{2}}b\ps{2}
\otimes\dots\otimes \wbar{h^{n-1}\ps{n}h^n\ps{n}}b\ps{n} \otimes
\wbar{h^n\ps{n+1}}b\ps{n+1} \otimes \wbar{1}\big) \otimes_{\cal H} h^0h^1\ps{1}\cdots h^n\ps{1}b\ps{1}
\\
\qquad
\stackrel{\eqref{coinviter}}{=}
\big(\wbar{h^1\ps{2}\cdots h^n\ps{2}}b\ps{2}
\otimes\dots\otimes \wbar{h^{n-1}\ps{n}h^n\ps{n}}b\ps{n} \otimes
\wbar{h^n\ps{n+1}}b\ps{n+1} \otimes \wbar{b\ps{n+2}}\big)\\
\qquad\hphantom{\stackrel{\eqref{coinviter}}{=}}{}
\otimes_{\cal H} h^0h^1\ps{1}\cdots h^n\ps{1}b\ps{1}
\\
\qquad{}
=\big(\wbar{h^1\ps{2}\cdots h^n\ps{2}}b\ps{2}
\otimes\dots\otimes \wbar{h^{n-1}\ps{n}h^n\ps{n}}b\ps{n} \otimes
\wbar{h^n\ps{n+1}}b\ps{n+1} \otimes \wbar{1}b\ps{n+2}\big)\\
\qquad{}
\hphantom{=}{}
 \otimes_{\cal H} h^0h^1\ps{1}\cdots h^n\ps{1}b\ps{1}
\\
\qquad
=\big(\wbar{h^1\ps{2}\cdots h^n\ps{2}}
\otimes\dots\otimes \wbar{h^{n-1}\ps{n}h^n\ps{n}} \otimes \wbar{h^n\ps{n+1}}
\otimes \wbar{1}\big)b\ps{2} \otimes_{\cal H} h^0h^1\ps{1}\cdots h^n\ps{1}b\ps{1}
\\
\qquad
=\big(\wbar{h^1\ps{2}\cdots h^n\ps{2}}
\otimes\dots\otimes \wbar{h^{n-1}\ps{n}h^n\ps{n}} \otimes \wbar{h^n\ps{n+1}}
\otimes \wbar{1}\big) \otimes_{\cal H} b\ps{2}\triangleright (h^0h^1\ps{1}\cdots h^n\ps{1}b\ps{1})
\\
\qquad{}
\stackrel{\eqref{HonH}}{=}
\big(\wbar{h^1\ps{2}\cdots h^n\ps{2}}
\otimes\dots\otimes \wbar{h^{n-1}\ps{n}h^n\ps{n}} \otimes \wbar{h^n\ps{n+1}}
\otimes \wbar{1}\big) \otimes_{\cal H} bh^0h^1\ps{1}\cdots h^n\ps{1}
\\
\qquad{}
=\varphi_n\big(\big[bh^0\otimes_{\cal B}\cdots\otimes_{\cal B} h^n\big]_{\cal B}\big),
\end{gather*}
that
\begin{gather*}
\varphi_n\big(\big[h^0\otimes_{\cal B} bh^{1}\otimes_{\cal B}\cdots\otimes_{\cal B} h^n\big]_{\cal B}\big)
\\
\qquad
=\big(\wbar{(bh^1)\ps{2}\cdots h^n\ps{2}} \otimes\cdots\otimes \wbar{h^{n-1}\ps{n}h^n\ps{n}} \otimes \wbar{h^n\ps{n+1}}
\otimes \wbar{1}\big) \otimes_{\cal H} h^0(bh^1)\ps{1}\cdots h^n\ps{1}
\\
\qquad
=\big(\wbar{b\ps{2}h^1\ps{2}\cdots h^n\ps{2}} \otimes\cdots\otimes \wbar{h^{n-1}\ps{n}h^n\ps{n}} \otimes
\wbar{h^n\ps{n+1}} \otimes \wbar{1}\big) \otimes_{\cal H} h^0b\ps{1}h^1\ps{1}\cdots h^n\ps{1}
\\
\qquad
=\big(\wbar{b\ps{2}}h^1\ps{2}\cdots h^n\ps{2} \otimes\cdots\otimes \wbar{h^{n-1}\ps{n}h^n\ps{n}} \otimes
\wbar{h^n\ps{n+1}} \otimes \wbar{1}\big) \otimes_{\cal H} h^0b\ps{1}h^1\ps{1}\cdots h^n\ps{1}
\\
\qquad
\stackrel{\eqref{coinv}}{=}
\big(\wbar{1}h^1\ps{2}\cdots h^n\ps{2} \otimes\cdots\otimes \wbar{h^{n-1}\ps{n}h^n\ps{n}} \otimes \wbar{h^n\ps{n+1}}
\otimes \wbar{1}\big) \otimes_{\cal H} h^0bh^1\ps{1}\cdots h^n\ps{1}
\\
\qquad
=\big(\wbar{h^1\ps{2}\cdots h^n\ps{2}} \otimes\cdots\otimes \wbar{h^{n-1}\ps{n}h^n\ps{n}} \otimes \wbar{h^n\ps{n+1}}
\otimes \wbar{1}\big) \otimes_{\cal H} h^0bh^1\ps{1}\cdots h^n\ps{1}
\\
\qquad
=\varphi_n\big(\big[h^0b\otimes_{\cal B} h^{1}\otimes_{\cal B}\cdots\otimes_{\cal B} h^n\big]_{\cal B}\big),
\end{gather*}
and that
\begin{gather*}
\varphi_n\big(\big[h^0\otimes_{\cal B}\cdots\otimes_{\cal B} h^i\otimes_{\cal B} bh^{i+1}\otimes_{\cal
B}\dots\otimes_{\cal B} h^n\big]_{\cal B}\big)
\\
\qquad
=\big(\wbar{h^1\ps{2}\cdots h^i\ps{2}(bh^{i+1})\ps{2}\cdots h^n\ps{2}}\otimes\dots\otimes
\wbar{h^i\ps{i+1}(bh^{i+1})\ps{i+1}\cdots h^n\ps{i+1}}
\\
\qquad\phantom{=}{}\otimes
\wbar{(bh^{i+1})\ps{i+2}\cdots h^n\ps{i+2}} \otimes \wbar{h^{i+2}\ps{i+3}\cdots
h^n\ps{i+3}}\otimes\cdots\otimes\wbar{h^n\ps{n+1}} \otimes \wbar{1}\big) \\
 \qquad\phantom{=}{}\otimes_{\cal H}
h^0h^1\ps{1}\cdots h^{i}\ps{1}\big(bh^{i+1}\big)\ps{1}\cdots h^n\ps{1}
\\
\qquad
=\big(\wbar{h^1\ps{2}\cdots h^i\ps{2}b\ps{2}h^{i+1}\ps{2}\cdots h^n\ps{2}}\otimes\dots\otimes
\wbar{h^i\ps{i+1}b\ps{i+1}h^{i+1}\ps{i+1}\cdots h^n\ps{i+1}}
\\
\qquad\phantom{=}{} \otimes
\wbar{b\ps{i+2}h^{i+1}\ps{i+2}\cdots h^n\ps{i+2}} \otimes \wbar{h^{i+2}\ps{i+3}\cdots
h^n\ps{i+3}}\otimes\cdots\otimes\wbar{h^n\ps{n+1}} \otimes \wbar{1}\big) \\
\qquad\phantom{=}{} \otimes_{\cal H}
h^0h^1\ps{1}\cdots h^{i}\ps{1}b\ps{1}h^{i+1}\ps{1}\cdots h^n\ps{1}
\\
\qquad
=\big(\wbar{h^1\ps{2}\cdots h^i\ps{2}b\ps{2}h^{i+1}\ps{2}\cdots h^n\ps{2}}\otimes\dots\otimes
\wbar{h^i\ps{i+1}b\ps{i+1}h^{i+1}\ps{i+1}\cdots h^n\ps{i+1}}
\\
\qquad\phantom{=}{} \otimes
\wbar{b\ps{i+2}}h^{i+1}\ps{i+2}\cdots h^n\ps{i+2} \otimes \wbar{h^{i+2}\ps{i+3}\cdots
h^n\ps{i+3}}\otimes\cdots\otimes\wbar{h^n\ps{n+1}} \otimes \wbar{1}\big) \\
\qquad\phantom{=}{}\otimes_{\cal H}
h^0h^1\ps{1}\cdots h^{i}\ps{1}b\ps{1}h^{i+1}\ps{1}\cdots h^n\ps{1}
\\
\qquad
\stackrel{\eqref{coinviter}}{=}
\big(\wbar{h^1\ps{2}\cdots h^i\ps{2}b\ps{2}h^{i+1}\ps{2}\cdots h^n\ps{2}}\otimes\dots\otimes
\wbar{h^i\ps{i+1}b\ps{i+1}h^{i+1}\ps{i+1}\cdots h^n\ps{i+1}}
\\
\qquad\phantom{=}{}
\otimes
\wbar{1}h^{i+1}\ps{i+2}\cdots h^n\ps{i+2} \otimes \wbar{h^{i+2}\ps{i+3}\cdots
h^n\ps{i+3}}\otimes\cdots\otimes\wbar{h^n\ps{n+1}} \otimes \wbar{1}\big) \\
\qquad\phantom{=}{} \otimes_{\cal H}
h^0h^1\ps{1}\cdots h^{i}\ps{1}b\ps{1}h^{i+1}\ps{1}\cdots h^n\ps{1}
\\
\qquad{}
=\big(\wbar{h^1\ps{2}\cdots (h^ib)\ps{2}h^{i+1}\ps{2}\cdots h^n\ps{2}}\otimes\cdots\otimes
\wbar{(h^ib)\ps{i+1}h^{i+1}\ps{i+1}\cdots h^n\ps{i+1}}
\\
\qquad\phantom{=}{} \otimes
\wbar{h^{i+1}\ps{i+2}\cdots h^n\ps{i+2}} \otimes \wbar{h^{i+2}\ps{i+3}\cdots
h^n\ps{i+3}}\otimes\cdots\otimes\wbar{h^n\ps{n+1}} \otimes \wbar{1}\big) \\
\qquad\phantom{=}{} \otimes_{\cal H}
h^0h^1\ps{1}\cdots (h^{i}b)\ps{1}h^{i+1}\ps{1}\cdots h^n\ps{1}
\\
\qquad{}
=\varphi_n\big(\big[h^0\otimes_{\cal B}\cdots\otimes_{\cal B} h^ib\otimes_{\cal B} h^{i+1}\otimes_{\cal
B}\cdots\otimes_{\cal B} h^n\big]_{\cal B}\big).
\end{gather*}

As for~\eqref{aux-map-psi}, we f\/irst note that by~\eqref{trans} the right hand side depends only on the representatives
appearing on the left hand side.
Moreover, for any $p\in{\cal H}$,
\begin{gather*}
\psi_n\big(\big(\wbar{g^0}\otimes\cdots\otimes \wbar{g^n}\big)\cdot p\otimes_{\cal H} h \big)
=\psi_n\big(\big(\wbar{g^0p\ps{1}}\otimes\cdots\otimes \wbar{g^np\ps{n+1}}\big)\otimes_{\cal H} h \big)
\\
\qquad
=\big[g^n\ps{2}p\ps{n+1}hS\big(g^0\ps{1}p\ps{1}\big) \otimes_{\cal B} g^0\ps{2}p\ps{2}S\big(g^1\ps{1}p\ps{3}\big) \otimes_{\cal B} \cdots\\
\qquad\phantom{=}{}
\otimes_{\cal B} g^{n-1}\ps{2}p\ps{2n}S\big(g^n\ps{1}p\ps{2n+1})\big]_{\cal B}
\\
\qquad
=\big[g^n\ps{2}p\ps{n+1}hS(p\ps{1})S\big(g^0\ps{1}\big) \otimes_{\cal B} g^0\ps{2}p\ps{2}S(p\ps{3})S\big(g^1\ps{1}\big) \otimes_{\cal B}
\cdots
\\
\qquad\phantom{=}{}
\otimes_{\cal B} g^{n-1}\ps{2}p\ps{2n}S(p\ps{2n+1})S\big(g^n\ps{1}\big)\big]_{\cal B}
\\
\qquad
=\big[g^n\ps{2}p\ps{2}hS(p\ps{1})S\big(g^0\ps{1}\big) \otimes_{\cal B} g^0\ps{2}S\big(g^1\ps{1}\big) \otimes_{\cal B} \cdots \otimes_{\cal B}
g^{n-1}\ps{2}S\big(g^n\ps{1}\big)\big]_{\cal B}
\\
\qquad
=\psi_n\big(\big(\wbar{g^0}\otimes\cdots\otimes \wbar{g^n}\big)\otimes_{\cal H} p\triangleright h \big).
\end{gather*}
Accordingly, $\psi_n$ is well-def\/ined for any $n\geq 0$.

Finally, we prove that $\varphi_n$ and $\psi_n$ are inverses to each other for any $n\geq 0$.
On one hand we have
\begin{gather*}
(\psi_n \circ \varphi_n)\big(\big[h^0\otimes_{\cal B}\cdots\otimes_{\cal B} h^n\big]_{\cal B} \big)
\\
\qquad
=\psi_n\big(\big(\wbar{h^1\ps{2}\cdots h^n\ps{2}} \otimes\cdots\otimes \wbar{h^{n-1}\ps{n}h^n\ps{n}} \otimes
\wbar{h^n\ps{n+1}} \otimes \wbar{1}\big) \otimes_{\cal H} h^0h^1\ps{1}\cdots h^n\ps{1} \big)
\\
\qquad
=\big[1\ps{2}h^0h^1\ps{1}\cdots h^n\ps{1}S\big(\big(h^1\ps{2}\cdots h^n\ps{2}\big)\ps{1}\big)\otimes_{\cal B} \big(h^1\ps{2}\cdots
h^n\ps{2}\big)\ps{2}S\big(\big(h^2\ps{3}\cdots h^n\ps{3}\big)\ps{1}\big)
\\
\qquad\phantom{=}{}
\otimes_{\cal B}\cdots \otimes_{\cal B} \big(h^{n-1}\ps{n}h^n\ps{n}\big)\ps{2}S\big(h^n\ps{n+1}\ps{1}\big)\otimes_{\cal B}
h^n\ps{n+1}\ps{2}S(1\ps{1})\big]_{\cal B}
\\
\qquad
=\big[h^0h^1\ps{1}\cdots h^n\ps{1}S\big(\big(h^1\ps{2}\cdots h^n\ps{2}\big)\ps{1}\big)\otimes_{\cal B} \big(h^1\ps{2}\cdots
h^n\ps{2}\big)\ps{2}S\big(\big(h^2\ps{3}\cdots h^n\ps{3}\big)\ps{1}\big)
\\
\qquad\phantom{=}{}
\otimes_{\cal B}\cdots \otimes_{\cal B} \big(h^{n-1}\ps{n}h^n\ps{n}\big)\ps{2}S\big(h^n\ps{n+1}\ps{1}\big)\otimes_{\cal B}
h^n\ps{n+1}\ps{2}\big]_{\cal B}
\\
\qquad
=\big[h^0h^1\ps{1}\cdots h^n\ps{1}S\big(h^1\ps{2}\ps{1}\cdots h^n\ps{2}\ps{1}\big)\otimes_{\cal B} h^1\ps{2}\ps{2}\cdots
h^n\ps{2}\ps{2}S\big(h^2\ps{3}\ps{1}\cdots h^n\ps{3}\ps{1}\big)
\\
\qquad\phantom{=}{}
\otimes_{\cal B}\cdots \otimes_{\cal B} h^{n-1}\ps{n}\ps{2}h^n\ps{n}\ps{2}S\big(h^n\ps{n+1}\ps{1}\big)\otimes_{\cal B}
h^n\ps{n+1}\ps{2}\big]_{\cal B}
\\
\qquad
=\big[h^0\cdot h^1\ps{1}\cdots h^n\ps{1}S\big(h^n\ps{2}\big)\cdots S\big(h^1\ps{2}\big)\otimes_{\cal B} h^1\ps{3}\cdot
h^2\ps{3}\cdots h^n\ps{3}S\big(h^n\ps{3}\big)\cdots S\big(h^2\ps{3}\big)
\\
\qquad\phantom{=}{}
\otimes_{\cal B}\cdots \otimes_{\cal B} h^{n-1}\ps{2n-1}h^n\ps{2n-1}S\big(h^n\ps{2n}\big)\otimes_{\cal B} h^n\ps{2n+1}\big]_{\cal B}
\\
\qquad
=\big[h^0\cdot h^1\ps{1}\cdots h^n\ps{1}S\big(h^n\ps{2}\big)\cdots S\big(h^1\ps{2}\big)\otimes_{\cal B} h^1\ps{3}\cdot
h^2\ps{3}\cdots h^n\ps{3}S\big(h^n\ps{3}\big)\cdots S\big(h^2\ps{3}\big)
\\
\qquad\phantom{=}{}
\otimes_{\cal B}\cdots \otimes_{\cal B} h^{n-1}\ps{2n-1}h^n\ps{2n-1}S\big(h^n\ps{2n}\big)\otimes_{\cal B} h^n\ps{2n+1}\big]_{\cal B}
\\
\qquad
=\big[h^0\cdot h^1\ps{1}\cdots h^{n-1}\ps{1}S\big(h^{n-1}\ps{2}\big)\cdots S\big(h^1\ps{2}\big)
\\
\qquad\phantom{=}{}
\otimes_{\cal B}
h^1\ps{3}\cdot h^2\ps{3}\cdots h^{n-1}\ps{3}S\big(h^{n-1}\ps{3}\big)\cdots S\big(h^2\ps{3}\big)\otimes_{\cal B}
\cdots \otimes_{\cal B} h^{n-1}\ps{2n-1}\otimes_{\cal B} h^n]_{\cal B} =\cdots
\\
\qquad{}
=\big[h^0\otimes_{\cal B}\cdots\otimes_{\cal B} h^n\big]_{\cal B},
\end{gather*}
while on the other hand
\begin{gather*}
(\varphi_n \circ \psi_n)\big((\wbar{g^0}\otimes\cdots\otimes \wbar{g^n})\otimes_{\cal H} h \big)
\\
\qquad
=\varphi_n\big([g^n\ps{2}hS(g^0\ps{1}) \otimes_{\cal B} g^0\ps{2}S(g^1\ps{1}) \otimes_{\cal B} \cdots \otimes_{\cal B}
g^{n-1}\ps{2}S(g^n\ps{1})]_{\cal B} \big)
\\
\qquad
=\big(\wbar{g^0\ps{2}\ps{2}S(g^1\ps{1})\ps{2}g^1\ps{2}\ps{2}S(g^2\ps{1})\ps{2} \cdots
g^{n-1}\ps{2}\ps{2}S(g^n\ps{1})\ps{2}} \otimes \cdots
\\
\qquad
\phantom{=}{}
\otimes \wbar{g^{n-2}\ps{2}\ps{n}S(g^{n-1}\ps{1})\ps{n}g^{n-1}\ps{2}\ps{n}S(g^n\ps{1})\ps{n}}\otimes
\wbar{g^{n-1}\ps{2}\ps{n+1}S(g^n\ps{1})\ps{n+1}} \otimes \wbar{1}\big)
\\
\qquad\phantom{=}{}\otimes_{\cal H}
g^n\ps{2}hS(g^0\ps{1})g^0\ps{2}\ps{1}S(g^1\ps{1})\ps{1}g^1\ps{2}\ps{1}S(g^2\ps{1})\ps{1}\cdots
g^{n-1}\ps{2}\ps{1}S(g^n\ps{1})\ps{1}
\\
\qquad
=\big(\wbar{g^0\ps{3}S(g^1\ps{1})g^1\ps{4}S(g^2\ps{2}) \cdots g^{n-1}\ps{n+2}S(g^n\ps{n})} \otimes \cdots
\\
\qquad
\phantom{=}
\otimes \wbar{g^{n-2}\ps{2n-1}S(g^{n-1}\ps{1})g^{n-1}\ps{2n}S(g^n\ps{2})}\otimes
\wbar{g^{n-1}\ps{2n+1}S(g^n\ps{1})} \otimes \wbar{1}\big)
\\
\qquad\phantom{=}\otimes_{\cal H}
g^n\ps{n+2}hS(g^0\ps{1})g^0\ps{2}S(g^1\ps{2})g^1\ps{3}S(g^2\ps{3})\cdots g^{n-1}\ps{n+1}S(g^n\ps{n+1})
\\
\qquad
=\big(\wbar{g^0S(g^n\ps{n})} \otimes \cdots\otimes \wbar{g^{n-2}S(g^n\ps{2})}\otimes \wbar{g^{n-1}S(g^n\ps{1})} \otimes
\wbar{1}\big)\otimes_{\cal H} g^n\ps{n+2}hS(g^n\ps{n+1})
\\
\qquad
=\big(\wbar{g^0S(g^n\ps{1}\ps{n})} \otimes\cdots\otimes \wbar{g^{n-1}S(g^n\ps{1}\ps{1})}\otimes
\wbar{1}\big)\otimes_{\cal H} g^n\ps{2}hS(g^n\ps{1}\ps{n+1})
\\
\qquad
=\big(\big(\wbar{g^0} \otimes\cdots\otimes \wbar{g^{n-1}}\big)\cdot S\big(g^n\ps{1}\big)\otimes \wbar{1}\big)\otimes_{\cal H}
g^n\ps{2} \triangleright h
\\
\qquad
=\big(\big(\wbar{g^0} \otimes\cdots\otimes \wbar{g^{n-1}}\big)\cdot S(g^n\ps{1})\otimes \wbar{1}\big)\cdot g^n\ps{2}
\otimes_{\cal H} h
\\
\qquad
=\big(\wbar{g^0} \otimes\cdots\otimes \wbar{g^{n-1}}\otimes \wbar{g^n}\big) \otimes_{\cal H} h. \tag*{\qed}
\end{gather*}
\renewcommand{\qed}{}
\end{proof}

We are now ready to state our main result.
\begin{Theorem}
\label{dirpict}
Let ${\cal I}\subseteq {\cal H}$ be a~coideal right ideal in a~Hopf algebra ${\cal H}$
such that ${\cal H}^{\co{\cal H}/{\cal I}}\subseteq {\cal H}$ is an ${\cal H}/{\cal I}$-Galois extension.
Let also $\ad({\cal H})={\cal H}$ be the left-right SAYD module with the right adjoint action, and the left
coaction given by the comultiplication of~${\cal H}$.
Then there exists an isomorphism
\begin{gather*}
\psi_{n}\colon \  {\rm C}_{n}({\cal H}/{\cal I},\ad({\cal H}))_{{\cal H}}\longrightarrow {\rm C}_{n}\big({\cal H}\mid{\cal H}^{\co{\cal H}/{\cal I}}\big)
\end{gather*}
of cyclic modules, defined by~\eqref{aux-map-psi}, \eqref{aux-map-vp}.
\end{Theorem}
\begin{proof}
Let us, as above, adopt the notation ${\cal B}:={\cal H}^{\co{\cal H}/{\cal I}}$.
We shall f\/irst go through the commutation with the face operators.
For $i=0$,
\begin{gather*}
\varphi_{n-1}d_0\big(\big[h^0\otimes_{\cal B}\cdots\otimes_{\cal B} h^n\big]_{\cal B}\big) =\varphi_{n-1}\big(\big[h^0h^1\otimes_{\cal
B}\cdots\otimes_{\cal B} h^n\big]_{\cal B}\big)
\\
\qquad
=\big(\wbar{h^2\ps{2}\cdots h^n\ps{2}} \otimes\cdots\otimes \wbar{h^{n-1}\ps{n-1}h^n\ps{n-1}} \otimes \wbar{h^n\ps{n}}
\otimes \wbar{1}\big) \otimes_{\cal H} h^0h^1h^2\ps{1}\cdots h^n\ps{1}
\\
\qquad
=d_0\varphi_{n}\big(\big[h^0\otimes_{\cal B}\cdots\otimes_{\cal B} h^n\big]_{\cal B}\big).
\end{gather*}
For $1\leq i \leq n-1$, we observe that
\begin{gather*}
\varphi_{n-1}d_i\big(\big[h^0\otimes_{\cal B}\cdots\otimes_{\cal B} h^n\big]_{\cal B}\big) =\varphi_{n-1}([h^0\otimes_{\cal
B}\cdots\otimes_{\cal B} h^ih^{i+1}\otimes_{\cal B}\cdots\otimes_{\cal B} h^n]_{\cal B})
\\
\qquad
=\big(\wbar{h^1\ps{2}\cdots(h^ih^{i+1})\ps{2}\cdots h^n\ps{2}} \otimes\cdots\otimes \wbar{(h^ih^{i+1})\ps{i+1}\dots
h^n\ps{i+1}}
\\
\qquad
\phantom{=}{}\otimes\wbar{h^{i+2}\ps{i+2}\cdots h^n\ps{i+2}}\otimes\cdots
\otimes\wbar{h^{n-1}\ps{n-1}h^n\ps{n-1}} \otimes \wbar{h^n\ps{n}} \otimes \wbar{1}\big)\\
\qquad
\phantom{=}{}
 \otimes_{\cal H}
h^0h^1\ps{1}\cdots(h^ih^{i+1})\ps{1}\cdots h^n\ps{1}
\\
\qquad
=\big(\wbar{h^1\ps{2}\cdots h^n\ps{2}} \otimes\cdots\otimes \wbar{h^i\ps{i+1}h^{i+1}\ps{i+1}\cdots h^n\ps{i+1}}\otimes
\varepsilon\big(\wbar{h^{i+1}\ps{i+2}\cdots h^n\ps{i+2}}\big) \otimes\cdots
\\
\qquad
\phantom{=}{}
\otimes \wbar{h^{n-1}\ps{n}h^n\ps{n}} \otimes \wbar{h^n\ps{n+1}} \otimes \wbar{1}\big) \otimes_{\cal H}
h^0h^1\ps{1}\cdots h^n\ps{1}
\\
\qquad
=d_i\varphi_{n}\big(\big[h^0\otimes_{\cal B}\cdots\otimes_{\cal B} h^n\big]_{\cal B}\big).
\end{gather*}
Finally for the last face operator we have
\begin{gather*}
\varphi_{n-1}d_n\big(\big[h^0\otimes_{\cal B}\cdots\otimes_{\cal B} h^n\big]_{\cal B}\big) =\varphi_{n-1}\big(\big[h^nh^0\otimes_{\cal
B}\cdots\otimes_{\cal B} h^{n-1}\big]_{\cal B}\big)
\\
\qquad
=\big(\wbar{h^1\ps{2}\cdots h^{n-1}\ps{2}} \otimes \cdots
\otimes \wbar{h^{n-2}\ps{n-1}h^{n-1}\ps{n-1}} \otimes \wbar{h^{n-1}\ps{n}} \otimes
\wbar{1}\big)
 \otimes_{\cal H} h^nh^0h^1\ps{1}\cdots h^{n-1}\ps{1}
\\
\qquad
\stackrel{\eqref{HonH}}{=}
\big(\wbar{h^1\ps{2}\cdots h^{n-1}\ps{2}} \otimes\cdots
\otimes \wbar{h^{n-2}\ps{n-1}h^{n-1}\ps{n-1}} \otimes \wbar{h^{n-1}\ps{n}} \otimes
\wbar{1}\big) \\
\qquad\hphantom{=}{}
\otimes_{\cal H} h^n\ps{2}\triangleright \big(h^0h^1\ps{1}\cdots h^{n-1}\ps{1}h^n\ps{1}\big)
\\
\qquad
=\big(\wbar{h^1\ps{2}\cdots h^{n-1}\ps{2}} \otimes\cdots
\otimes \wbar{h^{n-2}\ps{n-1}h^{n-1}\ps{n-1}} \otimes \wbar{h^{n-1}\ps{n}} \otimes \wbar{1}\big)\cdot
h^n\ps{2} \\
\qquad
\phantom{=}
\otimes_{\cal H} h^0h^1\ps{1}\cdots h^{n-1}\ps{1}h^n\ps{1}
\\
\qquad
=\big(\wbar{h^1\ps{2}\cdots h^{n}\ps{2}} \otimes\dots\otimes \wbar{h^{n-1}\ps{n}h^{n}\ps{n}} \otimes
\wbar{h^{n}\ps{n+1}} \big) \otimes_{\cal H} h^0h^1\ps{1}\cdots h^{n}\ps{1}
\\
\qquad
=\big(\wbar{h^1\ps{2}\cdots h^{n}\ps{2}} \otimes\cdots\otimes \wbar{h^{n-1}\ps{n}h^{n}\ps{n}} \otimes
\wbar{h^{n}\ps{n+1}}\otimes\varepsilon\big(\wbar{1}\big) \big) \otimes_{\cal H} h^0h^1\ps{1}\cdots h^{n}\ps{1}
\\
\qquad
=d_n\varphi_{n}\big(\big[h^0\otimes_{\cal B}\cdots\otimes_{\cal B} h^n\big]_{\cal B}\big).
\end{gather*}
We next investigate the interaction with the degeneracy operators.
To this end, for $0\leq j \leq n-1$ we observe that
\begin{gather*}
\varphi_{n+1}s_j\big(\big[h^0\otimes_{\cal B}\cdots\otimes_{\cal B} h^n\big]_{\cal B}\big)
=\varphi_{n+1}\big(\big[h^0\otimes_{\cal B}\cdots\otimes_{\cal B} h^j\otimes_{\cal B} 1\otimes_{\cal B} h^{j+1}\otimes_{\cal
B}\cdots\otimes_{\cal B} h^{n}\big]_{\cal B}\big)
\\
\qquad
=\big(\wbar{h^1\ps{2}\cdots h^n\ps{2}} \otimes\cdots
\otimes \wbar{h^j\ps{j+1}\cdots h^n\ps{j+1}} \otimes \wbar{h^{j+1}\ps{j+2}\cdots h^n\ps{j+2}}\\
\qquad
\phantom{=}{}
\otimes
\wbar{h^{j+1}\ps{j+3}\dots h^n\ps{j+3}}\otimes\cdots
\\
\qquad
\phantom{=}{}
\otimes \wbar{h^{n-1}\ps{n+1}h^n\ps{n+1}} \otimes \wbar{h^n\ps{n+2}} \otimes \wbar{1}\big) \otimes_{\cal H}
h^0h^1\ps{1}\cdots h^n\ps{1}
\\
\qquad
=\big(\wbar{h^1\ps{2}\cdots h^n\ps{2}} \otimes\cdots\otimes \wbar{h^j\ps{j+1}\cdots h^n\ps{j+1}} \otimes
\Delta\big(\wbar{h^{j+1}\ps{j+2}\cdots h^n\ps{j+2}}\big)\otimes \cdots
\\
\qquad
\phantom{=}{}
\otimes \wbar{h^{n-1}\ps{n}h^n\ps{n}} \otimes \wbar{h^n\ps{n+1}} \otimes \wbar{1}\big) \otimes_{\cal H}
h^0h^1\ps{1}\cdots h^n\ps{1}
\\
\qquad
=s_j\varphi_{n}\big(\big[h^0\otimes_{\cal B}\cdots\otimes_{\cal B} h^n\big]_{\cal B}\big).
\end{gather*}
Let us f\/inally check the cyclic operators.
To this end we have
\begin{gather*}
\varphi_{n}t_n\big(\big[h^0\otimes_{\cal B}\cdots\otimes_{\cal B} h^n\big]_{\cal B}\big) =\varphi_{n+1}\big(\big[h^n\otimes_{\cal B}
h^0\otimes_{\cal B}\cdots\otimes_{\cal B} h^{n-1}\big]_{\cal B}\big)
\\
\qquad
=\big(\wbar{h^0\ps{2}\cdots h^{n-1}\ps{2}} \otimes\cdots\otimes \wbar{h^{n-2}\ps{n}h^{n-1}\ps{n}} \otimes
\wbar{h^{n-1}\ps{n+1}} \otimes \wbar{1}\big) \otimes_{\cal H} h^nh^0\ps{1}\cdots h^{n-1}\ps{1}
\\
\qquad
=\big(\wbar{h^0\ps{2}h^1\ps{2}\cdots h^n\ps{2}}\otimes \wbar{h^1\ps{3}\cdots h^n\ps{3}} \otimes\cdots
\\
\qquad
\phantom{=}{}
\otimes \wbar{h^{n-1}\ps{n+1}h^n\ps{n+1}} \otimes \wbar{h^n\ps{n+2}} \big) \otimes_{\cal H}
h^0\ps{1}h^1\ps{1}\cdots h^n\ps{1}
\\
\qquad
=\big(\wbar{1}\cdot (h^0h^1\ps{1}\cdots h^n\ps{1})\ps{2}\otimes \wbar{h^1\ps{2}\cdots h^n\ps{2}} \otimes\cdots
\\
\qquad
\phantom{=}{}
\otimes \wbar{h^{n-1}\ps{n}h^n\ps{n}} \otimes \wbar{h^n\ps{n+1}} \big) \otimes_{\cal H}
(h^0h^1\ps{1}\cdots h^n\ps{1})\ps{1}
\\
\qquad
=t_n\varphi_{n}\big(\big[h^0\otimes_{\cal B}\cdots\otimes_{\cal B} h^n\big]_{\cal B}\big). \tag*{\qed}
\end{gather*}
\renewcommand{\qed}{}
\end{proof}
